\begin{description}
\item[(10.1)] The maximum duration occurs if Earth passes exactly along the
diameter of the Sun. Now consider the following figure, where Sun rays are
depicted as traveling to the right towards the hypothetical distant observer
(reason why they can be considered parallel):
% FIGURE NEEDED: hand-drawn sketch -- Earth's orbit (arc at right), the Sun
% (circle at left) and the shaded isosceles triangle spanning the transit
% chord, with the angle 2*alpha at the Sun and alpha marked on the parallel
% rays (2021_T10_Solution.pdf, page 1).

The arc along the circumference associated to the transit, $2\alpha$, can be
easily found looking at the shaded triangle:
\[ \alpha = \sin^{-1}\!\left(\frac{R_\odot}{R_{orb}}\right) \]
\hfill(2 points)

being $R_{orb}$ the orbital radius of the Earth. Now the time of the transit
can be found by considering the angular velocity of the Earth, for instance by
means of the following proportionality relations:
\[ \frac{t_{tr}}{T_{orb}} = \frac{2\alpha}{2\pi}
   = \frac{\sin^{-1}\!\left(\frac{R_\odot}{R_{orb}}\right)}{\pi} \]
\hfill(2 points)

Substituting with the known values for these quantities we get, i.e.,
\[ t_{tr} \sim 12.97\ \mathrm{h} = 12\mathrm{h}\,58\mathrm{m} \]
\hfill(1 point)

\item[(10.2)] First of all it must be said that the minimum orbital period is
obtained if we assume that the transit occurs along the diameter of the star.
In other cases, the planet would be crossing a shorter path in front of the
star during the same time, meaning that it would have a smaller angular
velocity and therefore a longer period.

That said, it means that we can resort to the same expression of 10.1, yet
this time we need 2 additional elements:
\begin{itemize}
\item Using the small angle approximation:
\[ \sin(\alpha) \sim \alpha \]
\hfill(1 point)
\item Invoking the full expression for the orbital period:
\[ T_{orb} = \frac{2\pi}{\sqrt{GM_*}}\,R_{orb}^{3/2} \]
\hfill(1 point)
\end{itemize}

Combining these results, we get:
\[ \frac{t_{tr}}{T_{orb}}
   = \frac{t_{tr}}{\frac{2\pi}{\sqrt{GM_*}}R_{orb}^{3/2}}
   = \frac{\frac{R_*}{R_{orb}}}{\pi} \]
\hfill(2 points)

And solving for $R_{orb}$:
\[ R_{orb} = \frac{GM_*\,t_{tr}^2}{2^2 R_*^2} \]
\hfill(2 points)

Finally putting this last result back into the expression for the orbital
period:
\[ T = \frac{\pi G M_*}{4R^3}\,t_{tr}^3 \]
\hfill(2 points)

At this point, just a numerical evaluation is needed, though a more elegant
solution can be achieved by means of scaling relations by noting that this
very same expression must be true in the case of the Earth--Sun system as seen
from far away. Having that 31 minutes is the 4\% of 12.94h:
\[ T = 365.25\cdot\frac{0.1}{0.1^3}\cdot 0.04^3
     \sim 2.3\ \textit{d\'ias} % sic: Spanish "days" left in the source
     \sim 2\mathrm{d}\,7\mathrm{h} \]
\hfill(2 points)

These values were inspired by the planetary system Trappist 1.
\end{description}
